% Generated draft: requires mathematical and source review.
\documentclass[11pt]{article}
\usepackage[margin=25mm]{geometry}
\usepackage{fontspec}
\setmainfont{Latin Modern Roman}
\usepackage{amsmath,amssymb,mathtools,unicode-math}
\setmathfont{Latin Modern Math}
\usepackage{enumitem,needspace,xurl,hyperref}
\setlength{\parindent}{0pt}
\setlength{\parskip}{5pt}
\setlength{\emergencystretch}{3em}
\newcommand{\N}{\mathbb{N}}
\newcommand{\Z}{\mathbb{Z}}
\newcommand{\Q}{\mathbb{Q}}
\newcommand{\R}{\mathbb{R}}
\newcommand{\C}{\mathbb{C}}
\newcommand{\F}{\mathbb{F}}
\newcommand{\A}{\mathbb{A}}
\newcommand{\PP}{\mathbb{P}}
\newcommand{\E}{\mathbb{E}}
\newcommand{\eps}{\varepsilon}
\newcommand{\dd}{\,\mathrm{d}}
\begin{document}
\section*{UnsolvedMath ID: 3000089}
\textbf{Source problem code:} catalogue entry
\subsection{Definitions and mathematical statement}
The source asks for a characterization of hypergraphs admitting an orientation with a prescribed symmetric local edge-connectivity requirement, and in particular an orientation that is $k$-edge-connected within a specified node subset.

One common finite directed-hypergraph convention starts with $\mathcal H=(V,\mathcal E)$, each hyperedge nonempty, and chooses a head $h(e)\in e$ for each $e\in\mathcal E$, with tail set $e\setminus\{h(e)\}$. For $X\subseteq V$, define
\[
\delta^-_{\vec{\mathcal H}}(X)=
\{e:h(e)\in X,\ e\setminus X\ne\varnothing\},
\qquad
\lambda_{\vec{\mathcal H}}(s,t)=
\min_{t\in X,\ s\notin X}|\delta^-_{\vec{\mathcal H}}(X)|.
\]
Under this convention, the general target for a specified symmetric requirement $r(s,t)=r(t,s)\in\Z_{\geq0}$ is to characterize those $\mathcal H$ for which there is a head selection satisfying
\[
\lambda_{\vec{\mathcal H}}(s,t)\geq r(s,t)
\quad\text{for every ordered pair }s\ne t.
\]
The subset target would take $S\subseteq V$, $k\in\Z_{\geq0}$, and require this inequality with right side $k$ for every distinct $s,t\in S$.

These displays specify a standard conditional meaning, not a confirmed source definition. Hypergraph orientation can instead mean a tail/head partition or another structure, and hyperpath or cut connectivity conventions differ. Likewise, ``within $S$'' may constrain only the endpoints while permitting intermediate vertices outside $S$, as in the displayed all-$V$ cuts, or may require an induced/orientation restriction on $S$ itself. The supplied record does not choose among these meanings. It asks for a necessary-and-sufficient characterization, without proposing one particular family of inequalities or an algorithmic complexity target.

The supplied summary reports well-developed graph orientation theory but only partial understanding of the hypergraph local-connectivity/subset extension. No precise partial theorem, hypotheses, or citation is provided.
\subsection{Short English statement}
The problem seeks necessary and sufficient conditions for orienting a hypergraph to meet symmetric directed connectivity demands, especially among a chosen set of vertices. The orientation and subset-connectivity conventions are missing; a common head-selection/cut model is stated conditionally.
\subsection{Sources}
\begin{itemize}
\item Local UnsolvedMath registry, permanent ID 3000089, exact source code catalogue entry, ``Well-balanced orientations of hypergraphs''; status partially solved, classification PARTIAL-PROGRESS.
\item Egres Open, \emph{List of open problems}, indexed page ``Well-balanced orientations of hypergraphs,'' \url{https://oldlemon.cs.elte.hu/egres/open/List_of_open_problems}, HTML extraction accessed July 29, 2026. The provenance says the archived index lists the item as open but is not a current-resolution authority. No separate problem-page URL or primary reference is supplied.
\item The supplied research summary reports established graph theory and partial hypergraph progress, especially for prescribed-subset local connectivity, without named authors, exact results, or authoritative citations.
\end{itemize}
\end{document}
